\section{Purpose}

This is a maintainer's companion to \texttt{paper-formulation.tex}, not a second copy of it. It assumes that document is open alongside this one and does not restate its physics; equation numbers below (\eqref{eq:route-b-multiplier}, etc.) refer to it directly. What belongs here instead: where every measured number quoted in the paper actually comes from, the history of the one place the paper's content changed direction (the wall condition), the rest of this project's record of claims made and later retracted, and the open engineering questions that are honest to state and were not resolved by writing the paper.

The working discipline this project has followed, worth stating once rather than repeating at each entry below: a claim is either derived, cited, or measured and reproducible from a named script under version control. ``An earlier revision argued X'' is not evidence for or against X; the numbers are.

\section{Provenance ledger}
\label{sec:ledger}

Every quantitative or structural claim in the paper that is not a direct algebraic consequence of an equation already on the page is backed by one of the following. All are run in CI on every push; none is a one-off calculation kept only in a notebook.

\begin{center}
\begin{tabular}{p{0.30\textwidth}p{0.32\textwidth}p{0.32\textwidth}}
\textbf{Claim} & \textbf{Script} & \textbf{What it checks} \\
\hline
Shifted-Legendre pressure basis has degree $N$, exactly, at every projection & \path{derivations/verify_legendre_pressure_basis.jl} & Symbolic (Rational\{BigInt\}) degree identity for $\tilde P_n(\psi)$ and every quadrature-evaluated integrand in \S\ref{subsec:projections} \\
Wall-condition frozen-slope derivation, route B algebra, well-definedness, self-adjointness, sign & \path{derivations/feasibility_pinned_contact_line.jl} \S1--5 & Analytic truncation-error rate ($M^{-3/2}$); the retracted finite-difference artifact, reproduced deliberately; the assembled operator's constraint identity, self-adjointness and positive-semidefiniteness (\S5) \\
Compliance operator: self-adjoint pairing, compactness, resolvable dimension, semidefiniteness & \path{derivations/audit_compliance_operator.jl}, \path{derivations/audit_weak_determinacy.jl} & Frobenius asymmetry under candidate pairings; singular-value decay and its convergence in $M,L$; eigenvalue sign of the symmetrised operator \\
Nested-closure conditioning (order unity, flat in $\delta$) vs.\ joint system ($\sim10^{18}$) & \path{derivations/audit_nested_closure.jl} & Condition number of the inner Galerkin matrix across $\delta=10^{-2}\dots10^{-6}$; the tangency/non-intersection/positivity crossing table \\
Tangency degeneracy at the poles; the acceleration-level closure's non-viability & \path{derivations/verify_accel_closure.jl} & Symbolic confirmation that $\partial_\theta C$ is odd and $2\pi$-periodic for any state, and of the parity argument underlying \eqref{eq:tangency-degenerate} \\
The global (non-piecewise) pressure representation's identity-theorem and rank-deficiency obstructions & \path{derivations/probe_global_basis_argmin.jl} & Rank and conditioning of the combined kinematic/zero-pressure system; monotonicity of its integrated residual in $\theta_c$ \\
Contact-time metric definitions and their 17\% spread & \path{src/postprocessing.jl}, \path{test/test_postprocessing.jl} & \texttt{primary\_contact\_time} (longest interval) vs.\ \texttt{contact\_time} (sum) vs.\ span; synthetic-phase test pinning all three \\
Validation against experiment (droplet trajectory, $\sim1.4\times$ the error-bar half-length) & \path{scripts/validate_experimental.jl}, \path{data/reference/experimental/} & Digitised trajectories with error bars; \texttt{PROVENANCE.md} in that directory records the digitisation itself \\
Sweep worker-count ablation (measured, not assumed from \texttt{Sys.CPU\_THREADS}) & \path{scripts/sweep.jl} & Per-case marginal memory footprint and the timing knee across $W=1,2,4,\dots$ \\
\end{tabular}
\end{center}

\section{The wall condition: what changed, and why}
\label{sec:wall-history}

This is the one place the paper's content changed direction mid-project, and it is worth recording in full because the reversal was not a matter of taste.

The first implementation was the Dirichlet basis of \S\ref{sec:wall}'s discussion (\texttt{wall=:pinned} in the code, never called ``route A'' or ``route B'' outside this document): $\eta(b)=0$ holds as an identity of the basis, with no constraint to enforce. It was adopted on the understanding that its only cost was a quantified, if severe, violation of no-flux at the wall. That understanding was incomplete. The real cost, measured directly, is that the basis destroys exact conservation of bath volume: $|\int_0^b\eta\,r\,dr|$ reaches $2.79$ (i.e.\ $17.5R^3$) over a reference impact (\texttt{We=1.0958, Bo=0.017, Oh=0.006, M=L=60, N=3, nq=40, b=6, h0=3}), against the droplet's own $4\pi/3\approx4.19R^3$ -- volume created from nothing at roughly four times the droplet's size. For a model whose physical content is the transfer of displaced volume into radiated capillary waves, this was disqualifying once noticed, and had gone unstated up to that point.

The Neumann-basis multiplier construction, \eqref{eq:route-b-multiplier}, was proposed as the fix and was, for one revision of this project, itself rejected on the grounds that the Neumann basis ``cannot represent a pinned meniscus with a nonzero wall slope, converging at only $O(1/M)$.'' That rejection rested on a diagnostic that was measuring its own finite-difference grid: it differenced a series fitted on a fixed 4000-point grid across the grid's own last interval, a location \emph{inside} the boundary layer of width $\sim b/k_M$ that the truncated series necessarily has at $r=b$. The number returned scaled as $\Delta r\,k_M$ -- a property of the grid, not of the series -- and refining the grid at fixed $M=640$ drove it from $-0.033$ toward $0$, while the true slope, computed analytically, is $-0.208$. \path{derivations/feasibility_pinned_contact_line.jl} \S1b reproduces this artifact deliberately, on purpose, so the failure mode stays visible to anyone who next touches this file. Computed correctly, the truncation error falls as $M^{-3/2}$ and the slope is recovered pointwise at every $r<b$; \S\ref{sec:wall}'s frozen-wall-slope result then removes the motivation for representing a time-varying wall slope at all.

Once the diagnostic was corrected, the multiplier construction was re-examined and adopted, with the operator-level verification recorded in the ledger above and now stated as settled fact in \S\ref{sec:wall}. It is implemented as \texttt{wall=:clamped}, applied at every step including free flight -- a detail whose omission was itself a bug caught during implementation: the first version of the code patched only the Newton-solve path and left the pinning constraint unenforced whenever the droplet was airborne, silently, for a residual of $0.17$ before the fix.

\emph{The question the reversal does not settle.} None of the above establishes that the wall condition matters at the bath radius $b=6$ used throughout the paper's reference impact. It was, for a period, described in this project's own notes as ``negligible at $b=6$'' on the strength of a measured $\Delta\mathrm{CoR}=-0.001$ between \texttt{:free} and \texttt{:clamped} at that radius. That description is retracted, for a reason independent of the multiplier construction itself: $b=6$ was inherited from the reference-impact configuration and was never itself shown to be in an asymptotic, wall-independent regime. The parent paper \citep{AlventosaEtAl2023} uses $b=25R$, justified by an explicit check that reflected waves do not reach the droplet; no equivalent check has been run here. A same-day sweep under \texttt{:free} alone (no pinning involved) found $\mathrm{CoR}$ moving by roughly $11\%$ between $b=6$ and $b=4$, which means $b=6$ was already flagged, by this project's own measurement, as possibly not asymptotic before any wall-condition comparison was made. Reporting the wall condition as unimportant at a radius that was itself never validated as safely large is not a supported claim, and \S\ref{sec:open} records it as open rather than settled.

\section{Corrections to earlier claims}
\label{sec:corrections}

The wall condition (\S\ref{sec:wall-history}) is the largest reversal, but not the only one. Recorded here, without repeating the specific numbers where the ledger in \S\ref{sec:ledger} already carries them:

\begin{itemize}
\item \emph{The contact-time metric.} A first-to-last \emph{span} was quoted as ``the contact time'' throughout an earlier revision of this project, in roughly a dozen places, before being noticed to be inflated by detachment chatter -- brief, spurious re-entries into contact immediately after the physical rebound, a stepper artifact of the \texttt{just\_left} guard rather than a second bounce. \texttt{primary\_contact\_time}, the duration of the \emph{longest} contiguous contact interval, is the corrected metric; a first attempt at the fix used the \emph{first} interval instead, which is usually but not always the longest, and returns a value three orders of magnitude too small on an under-resolved run that chatters from onset. Fixing this metric also exposed a second, load-bearing bug: \texttt{run\_simulation} had never returned per-level phase labels at all, which meant \texttt{contact\_time} and the coefficient of restitution were unreachable from a real simulation run and had only ever been exercised against synthetic test data.
\item \emph{The self-adjointness diagnosis, twice wrong in opposite directions.} One revision attributed the compliance operator's droplet-block asymmetry to a single factor of $x$ and claimed the operator was monotone because its symmetric part dominated in Frobenius norm; both claims were unsupported (the two candidate repair mechanisms are indistinguishable at zero deformation, where the measurement was taken, and norm dominance is not evidence of definiteness). The next revision over-corrected, testing each of the two actual mismatches (measure and displacement component) in isolation, finding neither sufficient alone, and concluding self-adjointness was unattainable. Repairing both together, \eqref{eq:b_l-selfadjoint}, in fact makes the operator self-adjoint to machine precision; the formulation now used in the paper.
\item \emph{The withdrawn resolvable-rank law.} A claimed scaling $\mathrm{rank}\,\mathcal A\approx0.75\,k_Mr_c+2$, attributing the resolvable pressure dimension to the bath truncation alone, was withdrawn: its supporting study set $M=L$ throughout, making the claimed collapse variable collinear with an untested droplet-side one, and capped its largest cases at a fixed quadrature count rather than the operator's own rank. Decoupling $M$ and $L$ shows both truncations contribute, with no single-parameter collapse. The paper accordingly measures the resolvable dimension from the assembled operator at each step rather than predicting it.
\item \emph{The deleted acceleration-level closure.} An earlier revision argued the position-level Galerkin closure could not be well conditioned, since the pressure enters $\eta,z_{cm},z_d$ only at $O(\delta^2)$, and proposed differentiating the contact condition twice in time to raise the pressure's entry to $O(1)$. The premise was false: a scalar factor cannot change a condition number, and the measured ill-conditioning came entirely from mixing pressure columns (uniformly $O(\delta^2)$) with a $\theta_c$ column that carries no such factor. Removing $\theta_c$ from the linear system, not raising the closure's differential order, is what the measured $10^{18}\to O(1)$ conditioning gain in the ledger reflects. The acceleration-level material, and the ``hidden constraint drift'' problem it introduced, were deleted in full.
\item \emph{The piecewise-pressure ``double standard''.} An earlier revision argued that a pressure representation global on $[-1,1]$ (with $p\equiv0$ off the patch imposed weakly) was structurally disqualified, since a nonzero polynomial cannot vanish identically on a sub-interval. That objection applies equally to every other weak condition in this closure, including the Galerkin condition \eqref{eq:galerkin} itself, and was withdrawn as a double standard. The piecewise representation is retained regardless, for the independent reason given in \S\ref{subsec:pressure-representation}: with $\theta_c$ fixed as an outer parameter, exact vanishing off the patch is free.
\item \emph{The tangency-degeneracy misreading.} The parity identity \eqref{eq:tangency-degenerate} was first read as showing tangency to be structurally unusable as a closing condition and was discarded outright. That was too strong: the degeneracy is confined to the poles, where it is the correct physical statement (no contact extent exists to constrain), and tangency is serviceable everywhere else in its selector role, \eqref{eq:tangency-selector}. \cite{AgueroEtAl2026} rely on the identical degeneracy in their own formulation without difficulty, for exactly this reason.
\item \emph{The nesting-circularity misdiagnosis.} An argument that computing $\hat c$ from a downstream, already-fixed $\theta_c$ is circular (since $p$'s support depends on $\theta_c$) is correct, but was mistakenly read as ruling out the opposite nesting used in \S\ref{subsec:contact-angle}, in which $\theta_c$ is supplied first as a parameter and $\hat c$ solved afterward. The two directions were conflated; only the first is circular.
\item \emph{A literature note, not a retraction of this project's own claims.} \cite{AlventosaEtAl2023}'s printed Fourier--Bessel normaliser, $2/(bJ_1(k_m))^2$, is the correct form for a Dirichlet (pinned) basis evaluated at the wrong argument ($k_m$ rather than $k_mb$), inconsistent with the free-wall no-flux condition their own model imposes. \path{src/bessel_moments.jl} independently carried an inverted version of the ratio table illustrating this point for one revision, which has been corrected; the two errors are unrelated.
\end{itemize}

\section{Open questions}
\label{sec:open}

Stated plainly, in order of how much they bear on the results in the paper:

\begin{itemize}
\item \emph{Whether the bath radius $b=6$ is wall-independent has not been established, under either wall condition.} \S\ref{sec:wall-history} records the retraction of the claim that it is. This is prior to, and separate from, whether \texttt{:free} and \texttt{:clamped} agree at that radius; both questions are open.
\item \emph{Convergence of the pressure moments $c_m,b_l,f$ in $M$, $L$, $N$, and the spectral-filter cutoff, jointly, is not established.} The pointwise pressure is not expected to converge, since the compliance operator is compact; whether its low-order moments converge is a separate question the paper does not resolve, and cutoff sensitivity has been measured to be the same size as the residual $N$-dependence at the settings tested.
\item \emph{The contact-radius peak leads DNS} ($\tau\approx0.89$ against $1.46$), and refining $N$ moves it \emph{away} from DNS rather than toward it, implicating the unresolved pressure profile rather than a resolution deficit. No experimental contact-radius measurement exists to adjudicate between this and a definition mismatch (DNS thresholds a trapped gas film; this model's $\theta_c$ bounds the kinematically matched region).
\item \emph{Detachment chatter is routed around, not eliminated.} The primary-interval metric (\S\ref{sec:ledger}) reports the correct headline number regardless, but the chatter itself -- brief spurious re-entries at the true rebound -- remains a stepper artifact of the free-flight/contact transition guard, not a resolved piece of physics.
\end{itemize}
